\renewcommand{\thesection}{\arabic{section}}\renewcommand{\theHsection}{chapter.\arabic{section}}\setcounter{section}{5}\section{常微分方程初步}
\noindent\begin{minipage}{\textwidth}\ProblemHeading{例题}{6.1}
求解方程 $\displaystyle\frac{dy}{dx}=-\frac{x}{y}$。
\end{minipage}\par\vfill
\noindent\begin{minipage}{\textwidth}\ProblemHeading{例题}{6.2}
求解方程 $(1+x^2)y\,dy+\sqrt{1-y^2}\,dx=0$。
\end{minipage}\par\vfill
\noindent\begin{minipage}{\textwidth}\ProblemHeading{例题}{6.3}
求解方程 $\displaystyle\frac{dy}{dx}=\frac{x+y}{x-y}$。
\end{minipage}\par\vfill
\noindent\begin{minipage}{\textwidth}\ProblemHeading{例题}{6.4}
求解方程 $x\,dy=(y+\sqrt{x^2+y^2})\,dx$。
\end{minipage}\par\vfill
\newpage
\noindent\begin{minipage}{\textwidth}\ProblemHeading{例题}{6.5}
求解方程 $\displaystyle\frac{dy}{dx}-\frac{y}{x}=-1$。
\end{minipage}\par\vfill
\noindent\begin{minipage}{\textwidth}\ProblemHeading{例题}{6.6}
求解方程 $(2y^2+y-x)\,dy-y\,dx=0$。
\end{minipage}\par\vfill
\noindent\begin{minipage}{\textwidth}\ProblemHeading{例题}{6.7}
求解方程 $\displaystyle\frac{dy}{dx}=xy+x^3y^3$。
\end{minipage}\par\vfill
\noindent\begin{minipage}{\textwidth}\ProblemHeading{例题}{6.8}
求解二阶方程 $xy^{\prime\prime}+y^{\prime}=4x$（$x\ne0$）。
\end{minipage}\par\vfill
\newpage
\noindent\begin{minipage}{\textwidth}\ProblemHeading{例题}{6.9}
求方程 $y^{\prime\prime}-e^{2y}=0$ 满足初值条件 $y(0)=0$、$y^{\prime}(0)=1$ 的特解。
\end{minipage}\par\vfill
\noindent\begin{minipage}{\textwidth}\ProblemHeading{习题}{6.10}
求解下列可分离变量的微分方程：
\SubQ{1}{$(1+x^2)\,dy=y\,dx$。}
\SubQ{2}{$y^{\prime}=e^{x-y}$。}
\SubQ{3}{$xy^{\prime}+y=y^2$。}
\end{minipage}\par\vfill
\noindent\begin{minipage}{\textwidth}\ProblemHeading{习题}{6.11}
求解下列可分离变量的微分方程：
\SubQ{1}{$yy^{\prime}=(1-2x)/y$。}
\end{minipage}\par\vfill
\noindent\begin{minipage}{\textwidth}\ProblemHeading{习题}{6.12}
求解下列微分方程：
\SubQ{1}{$y^{\prime}=y^2/x^2-2$。}
\SubQ{2}{$y^{\prime}=y/x+x/y$。}
\SubQ{3}{$\displaystyle\frac{dx}{x^2-xy+y^2}=\frac{dy}{2y^2-xy}$。}
\end{minipage}\par\vfill
\newpage
\noindent\begin{minipage}{\textwidth}\ProblemHeading{习题}{6.13}
求解下列微分方程：
\SubQ{1}{$(x^2+3y^2)\,dx-2xy\,dy=0$。}
\end{minipage}\par\vfill
\noindent\begin{minipage}{\textwidth}\ProblemHeading{习题}{6.14（校订）}
考虑方程 $\displaystyle\frac{dy}{dx}=f\!\left(\frac{a_1x+b_1y+c_1}{a_2x+b_2y+c_2}\right)$，其中系数为常数，表达式在所讨论区域有定义。
\SubQ{1}{证明：当 $a_1b_2-a_2b_1\ne0$ 时，可通过平移化为齐次方程；并讨论系数行列式为零时如何通过代换求解。}
\SubQ{2}{求解 $\displaystyle\frac{dy}{dx}=\frac{x+y+3}{x-y+1}$。}
\SubQ{3}{求解 $\displaystyle\frac{dy}{dx}=\frac{2x+4y+3}{x+2y+1}$。}
\end{minipage}\par\vfill
\noindent\begin{minipage}{\textwidth}\ProblemHeading{习题}{6.15}
求下列线性方程和 Bernoulli 方程的解：
\SubQ{1}{$(1+x^2)y^{\prime}-2xy=(1+x^2)^2$。}
\SubQ{2}{$y^{\prime}+\frac{1-2x}{x}=1$。}
\SubQ{3}{$y^{\prime}=\frac{y}{x+y^3}$。}
\end{minipage}\par\vfill
\noindent\begin{minipage}{\textwidth}\ProblemHeading{习题}{6.16}
求下列线性方程和 Bernoulli 方程的解：
\SubQ{1}{$y^{\prime}+y/x=y^2\ln x$。}
\SubQ{2}{$y^{\prime}=y\tan x+y^2\cos x$。}
\SubQ{3}{$y-y^{\prime}\cos x=y^2(1-\sin x)\cos x$。}
\end{minipage}\par\vfill
\newpage
\noindent\begin{minipage}{\textwidth}\ProblemHeading{习题}{6.17}
求下列方程满足初值条件的特解：
\SubQ{1}{$y^{\prime}=\frac{y}{x}\ln\frac{y}{x}$，$y(1)=1$。}
\SubQ{2}{$y^{\prime}+y/x=\sin x/x$，$y(\pi)=1$。}
\end{minipage}\par\vfill
\noindent\begin{minipage}{\textwidth}\ProblemHeading{习题}{6.18}
求解下列微分方程：
\SubQ{1}{$y^{\prime}+x=\sqrt{x^2+y}$。}
\SubQ{2}{$y^{\prime}=\cos(x-y)$。}
\SubQ{3}{$y^{\prime}-e^{x-y}+e^x=0$。}
\end{minipage}\par\vfill
\noindent\begin{minipage}{\textwidth}\ProblemHeading{习题}{6.19}
求解下列微分方程：
\SubQ{1}{$y^{\prime}\sin y+x\cos y+x=0$。}
\end{minipage}\par\vfill
\noindent\begin{minipage}{\textwidth}\ProblemHeading{习题}{6.20}
试用常数变易法导出 Bernoulli 方程 $y^{\prime}+P(x)y=Q(x)y^n$（$n\ne0,1$）的通解。
\end{minipage}\par\vfill
\newpage
\noindent\begin{minipage}{\textwidth}\ProblemHeading{习题}{6.21}
一条曲线过点 $(2,3)$，其在坐标轴间的任意切线段被切点平分。求这条曲线。
\end{minipage}\par\vfill
\noindent\begin{minipage}{\textwidth}\ProblemHeading{习题}{6.22}
设 $f$ 处处连续，且对任意 $x\in\mathbb R$ 有 $f(x)=\int_0^x f(t)\,dt$。求 $f(x)$。
\end{minipage}\par\vfill
\noindent\begin{minipage}{\textwidth}\ProblemHeading{习题}{6.23}
已知镭的衰变速率与镭的现存量成正比，比例常数为 $k>0$。设开始时镭的量为 $a$，求 $t$ 时刻镭的量 $x(t)$。
\end{minipage}\par\vfill
\noindent\begin{minipage}{\textwidth}\ProblemHeading{习题}{6.24}
一汽艇以 $10\,\mathrm{km/h}$ 的速度在静水上运动，发动机关闭 $20\,\mathrm{s}$ 后速度降至 $6\,\mathrm{km/h}$。设水对汽艇的阻力与汽艇速度成正比。求：
\SubQ{1}{发动机停止 $2\,\mathrm{min}$ 后汽艇的速度。}
\SubQ{2}{发动机停止后的第 $1\,\mathrm{min}$ 内汽艇所走的路程。}
\end{minipage}\par\vfill
\newpage
\noindent\begin{minipage}{\textwidth}\ProblemHeading{习题}{6.25}
求解下列二阶方程：
\SubQ{1}{$xy^{\prime\prime}=y^{\prime}$。}
\SubQ{2}{$y^{\prime\prime}=y^{\prime}/x+x$。}
\SubQ{3}{$y^{\prime\prime}=y^{\prime}+x$。}
\end{minipage}\par\vfill
\noindent\begin{minipage}{\textwidth}\ProblemHeading{习题}{6.26}
求解下列二阶方程：
\SubQ{1}{$y^{\prime\prime}+(y^{\prime})^2=2e^{-y}$。}
\end{minipage}\par\vfill
\noindent\begin{minipage}{\textwidth}\ProblemHeading{习题}{6.27（校订）}
对下列初值问题，讨论是否存在满足给定条件的经典解；存在时求出该解：
\SubQ{1}{$y^{\prime\prime}=y^{\prime}/x+x^2/y^{\prime}$，$y(1)=1$，$y^{\prime}(1)=0$。}
\SubQ{2}{$y^3y^{\prime\prime}=-1$，$y(1)=1$，$y^{\prime}(1)=0$。}
\end{minipage}\par\vfill
\noindent\begin{minipage}{\textwidth}\ProblemHeading{例题}{6.28}
求方程 $xy^{\prime\prime}-y^{\prime}=0$ 的通解。
\end{minipage}\par\vfill
\newpage
\noindent\begin{minipage}{\textwidth}\ProblemHeading{例题}{6.29}
求解方程 $y^{\prime\prime}-3y^{\prime}+2y=0$。
\end{minipage}\par\vfill
\noindent\begin{minipage}{\textwidth}\ProblemHeading{例题}{6.30}
求解方程 $y^{\prime\prime}+y^{\prime}+y=0$。
\end{minipage}\par\vfill
\noindent\begin{minipage}{\textwidth}\ProblemHeading{习题}{6.31}
在下列方程中，已知一个特解 $y_1$，求方程的通解：
\SubQ{1}{$y^{\prime\prime}+2y^{\prime}/x+y=0$，$y_1=\sin x/x$。}
\SubQ{2}{$y^{\prime\prime}\sin^2x=2y$，$y_1=\cot x$。}
\SubQ{3}{$(1-x^2)y^{\prime\prime}-2xy^{\prime}+2y=0$，$y_1=x$。}
\end{minipage}\par\vfill
\noindent\begin{minipage}{\textwidth}\ProblemHeading{习题}{6.32}
先用观察法求下列齐次方程的一个非零特解，再求其通解：
\SubQ{1}{$x^2y^{\prime\prime}-2xy^{\prime}+2y=0$，$x\ne0$。}
\SubQ{2}{$xy^{\prime\prime}-(1+x)y^{\prime}+y=0$，$x\ne0$。}
\end{minipage}\par\vfill
\newpage
\noindent\begin{minipage}{\textwidth}\ProblemHeading{习题}{6.33}
已知方程 $(1+x^2)y^{\prime\prime}+2xy^{\prime}-6x^2-2=0$ 的一个特解为 $y_1=x^2$。求该方程满足初值条件 $y(-1)=0$、$y^{\prime}(-1)=0$ 的特解。
\end{minipage}\par\vfill
\noindent\begin{minipage}{\textwidth}\ProblemHeading{习题}{6.34}
求下列常系数齐次方程的通解：
\SubQ{1}{$y^{\prime\prime}-2y^{\prime}-y=0$。}
\SubQ{2}{$y^{\prime\prime}+2y^{\prime}+2y=0$。}
\SubQ{3}{$y^{\prime\prime}+y^{\prime}-6y=0$。}
\end{minipage}\par\vfill
\noindent\begin{minipage}{\textwidth}\ProblemHeading{习题}{6.35}
求下列常系数非齐次方程的一个特解：
\SubQ{1}{$y^{\prime\prime}+y=2\sin(x/2)$。}
\SubQ{2}{$y^{\prime\prime}-6y^{\prime}+9y=(x+1)e^{2x}$。}
\end{minipage}\par\vfill
\noindent\begin{minipage}{\textwidth}\ProblemHeading{习题}{6.36}
验证下列函数组在实轴上的线性相关性：
\SubQ{1}{$1,x,x^2,\ldots,x^n$ 线性无关。}
\SubQ{2}{$1,\cos^2x,\sin^2x$ 线性相关。}
\end{minipage}\par\vfill
\newpage
\noindent\begin{minipage}{\textwidth}\ProblemHeading{习题}{6.37}
证明：区间 $I$ 上任何线性相关的两个可微函数 $y_1,y_2$，它们的 Wronski 行列式 $W=y_1y_2^{\prime}-y_1^{\prime}y_2$ 恒为零。
\end{minipage}\par\vfill
\noindent\begin{minipage}{\textwidth}\ProblemHeading{习题}{6.38}
证明下列函数在 $(0,2)$ 上线性无关，但其 Wronski 行列式 $W=y_1y_2^{\prime}-y_1^{\prime}y_2$ 恒为零：\[y_1(x)=\begin{cases}(x-1)^2,&0\le x\le1,\\0,&1<x\le2,\end{cases}\qquad y_2(x)=\begin{cases}0,&0\le x\le1,\\(x-1)^2,&1<x\le2.\end{cases}\]
\end{minipage}\par\vfill
\noindent\begin{minipage}{\textwidth}\ProblemHeading{习题}{6.39}
求下列方程的通解，其中 $x=x(t)$，撇号表示对 $t$ 求导：
\SubQ{1}{$x^{\prime\prime\prime}+3x^{\prime\prime}+3x^{\prime}+x=0$。}
\SubQ{2}{$x^{\prime\prime\prime}-2x^{\prime\prime}+x^{\prime}-2x=0$。}
\SubQ{3}{$x^{(4)}-8x^{\prime\prime}+18x=0$。}
\end{minipage}\par\vfill
\noindent\begin{minipage}{\textwidth}\ProblemHeading{习题}{6.40}
求下列方程的通解，其中 $x=x(t)$，撇号表示对 $t$ 求导：
\SubQ{1}{$x^{(4)}+2x^{\prime\prime}+x=0$。}
\end{minipage}\par\vfill
